\section{Lecciones aprendidas para nuestro PFC}

Nuestro PFC es el sistema de gestión de presustentaciones de la UTEQ: Angular, Spring Boot con JWT y PostgreSQL. Ahí los estudiantes suben su anteproyecto en PDF y los docentes registran sus evaluaciones. La tabla~\ref{tab:riesgos} muestra tres riesgos que podrían provocar un fallo parecido al de Facebook.

\begin{table}[H]
\centering
\footnotesize
\renewcommand{\arraystretch}{1.1}
\caption{Riesgos del PFC y medida preventiva que incorporaremos en TA-PFC-E1.}
\label{tab:riesgos}
\begin{tabularx}{\linewidth}{@{}>{\raggedright\arraybackslash}p{0.27\linewidth} >{\raggedright\arraybackslash}X >{\raggedright\arraybackslash}p{0.33\linewidth}@{}}
\toprule
\textbf{Riesgo} & \textbf{Parecido con Facebook} & \textbf{Medida en TA-PFC-E1} \\
\midrule
Aplicar en producción un script de \texttt{docs/migraciones} sin probarlo. & Como el comando de mantenimiento: un cambio de rutina tumba todo. Con \texttt{ddl-auto=validate}, si el esquema no coincide, el backend no arranca. & Probar cada migración en una copia de la base, hacer un respaldo con \texttt{pg\_dump} antes y tener el script para revertirla. \\
Un solo backend, una sola base de datos y los PDF en el disco del servidor. & Como el \emph{backbone}: todo depende de un componente común, y si cae, nadie sube anteproyectos ni registra notas. & Respaldo automático diario de la base y de \texttt{uploads/}, y un \emph{health check} con Spring Boot Actuator que avise de la caída. \\
El acceso de administración pasa por el mismo login JWT. & Como los ingenieros sin acceso remoto: si falla la autenticación o cambia \texttt{JWT\_SECRET}, ni el administrador puede entrar a corregir. & Un procedimiento de acceso de emergencia (SSH al servidor y a la base) documentado fuera del sistema y probado. \\
\bottomrule
\end{tabularx}
\end{table}
